% SPDX-License-Identifier: CC-BY-SA-4.0
% Author: Matthieu Perrin

\begingroup

\begin{frame}[fragile]{Impossibilté du consensus}

  \begin{block}{Définitions}
    Une configuration $C$ est :
    \begin{description}
    \item[$v$-monovalente] si seule $v$ peut être décidée après ce point
    \item[monovalente] s'il existe $v$ tel que $C$ est $v$-monovalente 
    \item[bivalente] si $C$ n'est pas monovalente
    \item[critique] si $C$ est bivalente et tous ses fils sont monovalents
    \end{description}
  \end{block}

  \pause  

  \begin{alertblock}{Théorème}
    Il n'existe pas d'algorithme de consensus\\
    qui n'utilise que des lectures et écritures sur des variables partagées.
  \end{alertblock}

  \begin{exampleblock}{Démonstration}
    \begin{description}
    \item[Introduction] Supposons (par l'absurde) qu'il existe un tel algorithme $A$ ;
    \item[Lemme 1] $A$ possède une configuration critique ;
    \item[Lemme 2] $A$ ne possède pas de configuration critique ;
    \item[Conclusion] absurde : $A$ n'existe pas !
    \end{description}
  \end{exampleblock}

  \footnoteref{M. Fischer, N. Lynch, M. Paterson. \textit{Impossibility of Distributed Consensus with One Faulty Process.} JACM (1985)}
  
\end{frame}

\endgroup
\endinput
